\section{Autodualité de l'unité interne de Planck sous inversion d'orientation}

Le sélecteur CT108 construit

\begin{equation}
\theta_{108}=\left(\frac{54}{\pi}\right)^2,
\qquad
\ell_P=\frac{54}{\pi}\ell_0,
\qquad
t_P=\frac{54}{\pi}t_0,
\label{eq:planck-length-time}
\end{equation}

et

\begin{equation}
E_P=\frac{\hbar_{\rm int}}{t_P}
=\frac{\pi\hbar_{\rm int}}{54t_0}.
\label{eq:planck-energy}
\end{equation}

En comparant \cref{eq:boltzmann-clock,eq:planck-energy}, on obtient

\begin{equation}
\boxed{E_P=k_B\Theta_{\rm clk}\log3.}
\label{eq:planck-trit}
\end{equation}

La relation ne dépend pas d'une coïncidence entre unités : les deux membres
proviennent de la même structure chronologique \(108\).

Une fois la constante gravitationnelle dérivée en interne dans le
\cref{ch:gravity}, la famille de Planck satisfait

\begin{align}
\ell_P^2&=\frac{G\hbar}{c^3}=\theta_{108}\ell_0^2,
\label{eq:planck-area}\\
F_P&=\frac{c^4}{G},
\qquad
P_P=\frac{c^5}{G},
\label{eq:planck-force-power}\\
\rho_P&=\frac{\hbar}{\theta_{108}^2c^5t_0^4}.
\label{eq:planck-density}
\end{align}

Bien que \(G\) et \(\hbar\) se transforment réciproquement entre les sections
antérieure et postérieure au retour, leur produit dans \(\ell_P^2\) conserve la
géométrie. Cette invariance détermine l'aire comme interface entre les
réalisations géométrique, informationnelle et d'intrication.
